% Relativistic Chapter VII, §69--70: Viscous Couette Flow and Perturbation Equations
% Relativistic generalization of Chandrasekhar, Chapter VII §69--70
% Conventions: see RELATIVISTIC_CONVENTIONS.md

\section{Relativistic viscous Couette flow}\label{sec:rel-7-69}

In \S\ref{sec:7-69} we described the classical viscous Couette flow
between two coaxial rotating cylinders.  We now reformulate this
problem in the relativistic framework, replacing the Navier--Stokes
viscous stress by the Israel--Stewart causal dissipative theory.

\subsection{Relativistic energy-momentum tensor with Israel--Stewart
  viscosity}\label{sec:rel-7-69a}

The energy-momentum tensor for a relativistic dissipative fluid is
(cf.\ \textsc{RELATIVISTIC\_CONVENTIONS.md})
\begin{equation}
\label{eq:rel-7-134}
T^{\mu\nu} = \frac{\varepsilon + p}{c^{2}}\,u^{\mu}\,u^{\nu}
  + (p + \Pi)\,\Delta^{\mu\nu}
  + \pi^{\mu\nu}
  + \frac{1}{c^{2}}\bigl(q^{\mu}u^{\nu}+q^{\nu}u^{\mu}\bigr),
\end{equation}
where $\varepsilon$ is the energy density, $p$ the equilibrium
pressure, $\Pi$ the bulk viscous pressure, $\pi^{\mu\nu}$ the shear
stress tensor, and $q^{\mu}$ the heat flux four-vector.  The
projection tensor is
$\Delta^{\mu\nu}=g^{\mu\nu}+u^{\mu}u^{\nu}/c^{2}$ and the
four-velocity satisfies $u^{\mu}u_{\mu}=-c^{2}$.

In the Israel--Stewart theory the shear stress is not given
algebraically by the velocity gradient; instead it satisfies the
\emph{relaxation equation}
\begin{equation}
\label{eq:rel-7-135}
\tau_{\pi}\,u^{\alpha}\nabla_{\alpha}\pi^{\langle\mu\nu\rangle}
  + \pi^{\mu\nu}
  = 2\,\eta_{\mathrm{s}}\,\sigma^{\mu\nu},
\end{equation}
where $\tau_{\pi}>0$ is the shear relaxation time,
$\eta_{\mathrm{s}}$ is the coefficient of shear viscosity, and
$\sigma^{\mu\nu}$ is the relativistic shear tensor
\begin{equation}
\label{eq:rel-7-136}
\sigma^{\mu\nu}
  = \tfrac{1}{2}\bigl(\Delta^{\mu\alpha}\nabla_{\alpha}u^{\nu}
    + \Delta^{\nu\alpha}\nabla_{\alpha}u^{\mu}\bigr)
  - \tfrac{1}{3}\,\Delta^{\mu\nu}\,\nabla_{\alpha}u^{\alpha}.
\end{equation}
The angular brackets in~\eqref{eq:rel-7-135} denote the symmetric,
traceless, and transverse (to~$u^{\mu}$) projection.  The crucial
feature of~\eqref{eq:rel-7-135} is that $\tau_{\pi}>0$ renders the
viscous sector \emph{hyperbolic}: perturbations in $\pi^{\mu\nu}$
propagate at a finite speed
\begin{equation}
\label{eq:rel-7-137}
v_{\mathrm{visc}} = \sqrt{\frac{\eta_{\mathrm{s}}}
  {(\varepsilon+p)\,\tau_{\pi}/c^{2}}} \le c,
\end{equation}
so that causality is preserved.  When $\tau_{\pi}\to 0$
equation~\eqref{eq:rel-7-135} reduces to the Navier--Stokes
constitutive relation $\pi^{\mu\nu}=2\eta_{\mathrm{s}}\sigma^{\mu\nu}$,
which is acausal (parabolic).

Similarly, the bulk viscous pressure satisfies
\begin{equation}
\label{eq:rel-7-138}
\tau_{\Pi}\,u^{\alpha}\nabla_{\alpha}\Pi + \Pi
  = -\zeta\,\nabla_{\mu}u^{\mu},
\end{equation}
with relaxation time $\tau_{\Pi}>0$ and bulk viscosity coefficient
$\zeta$.  For incompressible flow ($\nabla_{\mu}u^{\mu}=0$) the bulk
viscous pressure vanishes in equilibrium, and we shall henceforth
set $\Pi=0$ and $q^{\mu}=0$ (no heat flux).

\subsection{Equilibrium Couette profile: relativistic
  corrections}\label{sec:rel-7-69b}

We work in cylindrical coordinates $(t,r,\varphi,z)$ on a flat
(Minkowski) spacetime.  The equilibrium four-velocity for purely
azimuthal flow with angular velocity $\Omega(r)$ is
\begin{equation}
\label{eq:rel-7-139}
u^{\mu} = \gamma\,(c,\,0,\,r\Omega,\,0),
\qquad
\gamma = \Bigl(1-\frac{r^{2}\Omega^{2}}{c^{2}}\Bigr)^{-1/2},
\end{equation}
where $\gamma$ is the Lorentz factor.  The azimuthal velocity is
$V(r)=r\Omega(r)$ and we require $V<c$ everywhere.

For the equilibrium state the shear tensor~\eqref{eq:rel-7-136}
reduces, after projection onto the azimuthal--radial plane, to the
single independent component
\begin{equation}
\label{eq:rel-7-140}
\sigma^{r\varphi}
  = \frac{\gamma^{2}}{2}\,r\,\frac{d\Omega}{dr}.
\end{equation}
In the Israel--Stewart equilibrium ($u^{\alpha}\nabla_{\alpha}
\pi^{\mu\nu}=0$), the relaxation equation~\eqref{eq:rel-7-135}
yields
\begin{equation}
\label{eq:rel-7-141}
\pi^{r\varphi}_{\mathrm{eq}}
  = 2\eta_{\mathrm{s}}\,\sigma^{r\varphi}
  = \eta_{\mathrm{s}}\,\gamma^{2}\,r\,\frac{d\Omega}{dr}.
\end{equation}

The radial component of the relativistic momentum equation
$\nabla_{\mu}T^{\mu\nu}=0$ yields the centrifugal balance
\begin{equation}
\label{eq:rel-7-142}
\frac{d}{dr}\!\Bigl(\frac{p}{\varepsilon+p}\Bigr)
  = \frac{\gamma^{2}\,V^{2}}{c^{2}\,r}
  + \frac{1}{\varepsilon+p}\,
    \frac{1}{r}\frac{d}{dr}\bigl(r\,\pi^{r\varphi}_{\mathrm{eq}}\bigr)
    - \text{curvature corrections}.
\end{equation}
The azimuthal component gives the equation for the equilibrium
velocity profile.  Defining the kinematic viscosity
$\nu=\eta_{\mathrm{s}}/w$ with
$w=(\varepsilon+p)/c^{2}$ the relativistic enthalpy density, the
azimuthal equilibrium condition becomes
\begin{equation}
\label{eq:rel-7-143}
\nu\,\gamma^{2}\,\frac{d}{dr}\!\left[
  \gamma^{2}\,\frac{d}{dr}\!\Bigl(
    \frac{\gamma\,V}{r}\Bigr)\,r^{3}
\right] = 0.
\end{equation}
In the non-relativistic limit $\gamma\to 1$ this reduces to
\begin{equation}
\label{eq:rel-7-144}
\nu\,\frac{d}{dr}\!\left(\frac{dV}{dr}+\frac{V}{r}\right) = 0,
\end{equation}
which is precisely equation~\eqref{eq:7-142} of \S\ref{sec:7-69},
confirming the classical Couette solution
$V=Ar+B/r$.

For the relativistic case, equation~\eqref{eq:rel-7-143} admits the
solution
\begin{equation}
\label{eq:rel-7-145}
V(r) = Ar + \frac{B}{r}
  + \frac{A^{3}r^{3}+3A^{2}Br+3AB^{2}/r+B^{3}/r^{3}}{2c^{2}}
  + O(V^{4}/c^{4}),
\end{equation}
so that the angular velocity acquires relativistic corrections:
\begin{equation}
\label{eq:rel-7-146}
\Omega(r) = A + \frac{B}{r^{2}}
  + \frac{1}{2c^{2}}\Bigl(A + \frac{B}{r^{2}}\Bigr)^{3}r^{2}
  + O(c^{-4}),
\end{equation}
where $A$ and $B$ are determined by the boundary conditions on the
two cylinders exactly as in the classical case
(equation~\eqref{eq:7-145}), with $O(V^{2}/c^{2})$ corrections.

The corresponding expression for the constants is
\begin{equation}
\label{eq:rel-7-147}
A = -\Omega_{1}\eta^{2}\,\frac{1-\mu/\eta^{2}}{1-\eta^{2}}
  \bigl[1+O(R_{1}^{2}\Omega_{1}^{2}/c^{2})\bigr],
\qquad
B = \Omega_{1}\,\frac{R_{1}^{2}(1-\mu)}{1-\eta^{2}}
  \bigl[1+O(R_{1}^{2}\Omega_{1}^{2}/c^{2})\bigr],
\end{equation}
with $\mu=\Omega_{2}/\Omega_{1}$ and $\eta=R_{1}/R_{2}$ retaining
their classical definitions (equation~\eqref{eq:7-146}).

\medskip
\noindent\textbf{Remark (causality).}  In the Navier--Stokes
formulation the viscous equilibrium~\eqref{eq:rel-7-144} is
instantaneously established.  In the Israel--Stewart theory, the
equilibrium is reached on a timescale $\sim\tau_{\pi}$; but once
$t\gg\tau_{\pi}$ the stationary profiles~\eqref{eq:rel-7-145}
and~\eqref{eq:rel-7-146} are the same.  The difference becomes
crucial only for the perturbation dynamics.

%% ================================================================
\section{The perturbation equations: relativistic
  formulation}\label{sec:rel-7-70}

We now investigate the stability of the relativistic Couette flow
described in \S\ref{sec:rel-7-69}.  The perturbation analysis follows
the structure of the classical treatment in \S\ref{sec:7-70}, but the
viscous sector is governed by the Israel--Stewart relaxation
equation~\eqref{eq:rel-7-135} rather than by the Navier--Stokes
stress.

\subsection{Linearized relativistic equations}\label{sec:rel-7-70a}

Let the perturbed state be characterised by
\begin{equation}
\label{eq:rel-7-148}
u^{\mu} = u^{\mu}_{\mathrm{eq}} + \delta u^{\mu},
\quad
p = p_{\mathrm{eq}} + \delta p,
\quad
\varepsilon = \varepsilon_{\mathrm{eq}} + \delta\varepsilon,
\quad
\pi^{\mu\nu} = \pi^{\mu\nu}_{\mathrm{eq}} + \delta\pi^{\mu\nu},
\end{equation}
where perturbations are taken to be axisymmetric
($\partial_{\varphi}=0$) and of first order.

Writing the perturbation velocity components in the local rest frame
as $u_{r}$, $u_{\varphi}$, $u_{z}$, and the pressure perturbation as
$\varpi=\delta p/w$ (with $w=(\varepsilon+p)/c^{2}$), the linearized
radial, azimuthal, and axial momentum equations become
\begin{align}
\label{eq:rel-7-149}
\gamma^{2}\frac{\partial u_{r}}{\partial t}
  - \frac{2\gamma^{2}V}{r}\,u_{\varphi}
  &= -\frac{\partial\varpi}{\partial r}
  + \frac{1}{w}\,\nabla_{\mu}\delta\pi^{\mu r}
  + \frac{\gamma^{2}V^{2}}{c^{2}r}\,
    \frac{\delta\varepsilon}{\varepsilon+p}\,,
\\[6pt]
\label{eq:rel-7-150}
\gamma^{2}\frac{\partial u_{\varphi}}{\partial t}
  + \gamma^{2}\!\left(\frac{dV}{dr}+\frac{V}{r}\right)u_{r}
  &= \frac{1}{w}\,\nabla_{\mu}\delta\pi^{\mu\varphi},
\\[6pt]
\label{eq:rel-7-151}
\frac{\partial u_{z}}{\partial t}
  &= -\frac{\partial\varpi}{\partial z}
  + \frac{1}{w}\,\nabla_{\mu}\delta\pi^{\mu z}.
\end{align}
The continuity equation (baryon conservation) for axisymmetric
perturbations is
\begin{equation}
\label{eq:rel-7-152}
\frac{\partial(\gamma\,u_{r})}{\partial r}
  + \frac{\gamma\,u_{r}}{r}
  + \frac{\partial u_{z}}{\partial z}
  + \frac{\gamma^{3}V}{c^{2}}\,\frac{\partial u_{r}}{\partial t}
  = 0.
\end{equation}
In the non-relativistic limit $\gamma\to 1$, the last term vanishes
and equations~\eqref{eq:rel-7-149}--\eqref{eq:rel-7-152} reduce
exactly to equations~\eqref{eq:7-149}--\eqref{eq:7-153} of
\S\ref{sec:7-70}.

\subsection{Israel--Stewart perturbation of the shear
  stress}\label{sec:rel-7-70b}

The key difference from the classical analysis is that the viscous
stress perturbation $\delta\pi^{\mu\nu}$ is not slaved to the
velocity gradient; it satisfies the linearised Israel--Stewart
equation
\begin{equation}
\label{eq:rel-7-153}
\tau_{\pi}\,\frac{\partial\,\delta\pi^{\mu\nu}}{\partial t}
  + \delta\pi^{\mu\nu}
  = 2\eta_{\mathrm{s}}\,\delta\sigma^{\mu\nu}
  - \tau_{\pi}\,(\delta u^{\alpha}\nabla_{\alpha})
    \pi^{\mu\nu}_{\mathrm{eq}},
\end{equation}
where $\delta\sigma^{\mu\nu}$ is the perturbed shear tensor.  In the
limit $\tau_{\pi}\to 0$ the left-hand side reduces to
$\delta\pi^{\mu\nu}=2\eta_{\mathrm{s}}\delta\sigma^{\mu\nu}$,
recovering the Navier--Stokes constitutive law and hence the
classical perturbation equations.

The explicit components of $\delta\sigma^{\mu\nu}$ in cylindrical
coordinates for axisymmetric perturbations are
\begin{align}
\label{eq:rel-7-154}
\delta\sigma^{rr} &= \gamma^{2}\,\frac{\partial u_{r}}{\partial r}
  - \tfrac{1}{3}\gamma^{2}\,\nabla_{\mu}\delta u^{\mu},
\\
\delta\sigma^{r\varphi} &= \tfrac{1}{2}\gamma^{2}
  \Bigl(\frac{\partial u_{\varphi}}{\partial r}
  - \frac{u_{\varphi}}{r}
  + \frac{\gamma^{2}V}{c^{2}}\frac{\partial u_{r}}{\partial t}\Bigr),
\\
\delta\sigma^{zz} &= \frac{\partial u_{z}}{\partial z}
  - \tfrac{1}{3}\gamma^{2}\,\nabla_{\mu}\delta u^{\mu}.
\end{align}

\subsection{Normal mode analysis}\label{sec:rel-7-70c}

Following the classical treatment, we decompose the perturbations
into normal modes:
\begin{equation}
\label{eq:rel-7-155}
\begin{aligned}
u_{r} &= e^{pt}\,u(r)\cos kz, &\quad
u_{\varphi} &= e^{pt}\,v(r)\sin kz, \\
u_{z} &= e^{pt}\,w(r)\cos kz, &\quad
\varpi &= e^{pt}\,\hat\varpi(r)\cos kz, \\
\delta\pi^{\mu\nu} &= e^{pt}\,\Pi^{\mu\nu}(r)
  \times\{\cos kz \text{ or } \sin kz\},
\end{aligned}
\end{equation}
where $k$ is the axial wave number and $p$ may be complex.  For the
Israel--Stewart stress components the relaxation
equation~\eqref{eq:rel-7-153} yields the algebraic relation (for each
mode)
\begin{equation}
\label{eq:rel-7-156}
\Pi^{\mu\nu}(r)
  = \frac{2\eta_{\mathrm{s}}}{1+\tau_{\pi}p}\,
    \delta\sigma^{\mu\nu}_{\mathrm{mode}}(r)
  - \frac{\tau_{\pi}}{1+\tau_{\pi}p}\,
    (\delta u^{\alpha}\nabla_{\alpha})
    \pi^{\mu\nu}_{\mathrm{eq}},
\end{equation}
so that the Israel--Stewart viscous term in the momentum equations
enters through an \emph{effective kinematic viscosity}
\begin{equation}
\label{eq:rel-7-157}
\nu_{\mathrm{eff}}(p) = \frac{\nu}{1+\tau_{\pi}p},
\end{equation}
which is frequency-dependent.  For steady perturbations ($p=0$),
$\nu_{\mathrm{eff}}=\nu$ and the classical eigenvalue problem is
recovered.

Substituting~\eqref{eq:rel-7-156} into the linearised momentum
equations and carrying out the same elimination of $w$ and $\varpi$
as in the classical \S\ref{sec:7-70}, we arrive at the pair of
coupled equations (measuring $r$ in units of $R_{2}$, with
$a^{2}=k^{2}R_{2}^{2}$ and
$\sigma=pR_{2}^{2}/\nu$):
\begin{equation}
\label{eq:rel-7-158}
\frac{1}{1+\tau_{\pi}p}\bigl(DD_{*}-a^{2}\bigr)
\Bigl[\frac{\gamma^{2}}{1+\tau_{\pi}p}\bigl(DD_{*}-a^{2}\bigr)
  - \sigma\gamma^{2}\Bigr]u
= -T_{\mathrm{rel}}\,a^{2}
  \Bigl(\frac{1}{r^{2}}-\kappa_{\mathrm{rel}}\Bigr)v,
\end{equation}
\begin{equation}
\label{eq:rel-7-159}
\frac{\gamma^{2}}{1+\tau_{\pi}p}
  \bigl(DD_{*}-a^{2}\bigr)v
  - \sigma\gamma^{2}v
= \bigl(D_{*}V\bigr)_{\mathrm{rel}}\,u,
\end{equation}
where $D$ and $D_{*}$ retain their classical meanings
(equation~\eqref{eq:7-106}) and we have introduced the
\emph{relativistic Taylor number}
\begin{equation}
\label{eq:rel-7-160}
T_{\mathrm{rel}}
  = -\frac{4A_{\mathrm{rel}}B_{\mathrm{rel}}}{R_{2}^{4}}
  = T_{\mathrm{cl}}
    \times\gamma_{1}^{2}\gamma_{2}^{2}
    \bigl[1+O(V^{2}/c^{2})\bigr],
\end{equation}
with $T_{\mathrm{cl}}$ the classical Taylor number of
equation~\eqref{eq:7-169}, and the relativistic analogue of
$\kappa$:
\begin{equation}
\label{eq:rel-7-161}
\kappa_{\mathrm{rel}}
  = -\frac{A_{\mathrm{rel}}\,R_{2}^{4}}{B_{\mathrm{rel}}}
  = \kappa\bigl[1+O(V^{2}/c^{2})\bigr].
\end{equation}
The notation $\gamma_{1}=\gamma(R_{1})$,
$\gamma_{2}=\gamma(R_{2})$ refers to the Lorentz factor evaluated
at each cylinder wall.

\medskip
\noindent\textbf{Remark.}  The structure of
equations~\eqref{eq:rel-7-158} and~\eqref{eq:rel-7-159} is
identical to the classical system~\eqref{eq:7-167}--\eqref{eq:7-168}
except for two modifications: (i) every occurrence of $\nu$ is
replaced by $\nu_{\mathrm{eff}}(p)=\nu/(1+\tau_{\pi}p)$, and
(ii)~factors of $\gamma^{2}$ multiply the inertial terms.

\subsection{Stability for $\mu>\eta^{2}$: relativistic
  generalisation}\label{sec:rel-7-70d}

We now establish that the Rayleigh criterion $\mu>\eta^{2}$ remains a
\emph{sufficient condition for stability} in the relativistic theory.

The argument follows the classical energy-integral method of
\S\ref{sec:7-70}.  Multiply equation~\eqref{eq:rel-7-158} by
$r\,u^{*}$ and equation~\eqref{eq:rel-7-159} by
$r\,T_{\mathrm{rel}}\,a^{2}\phi_{\mathrm{rel}}(r)\,v^{*}$, where
\begin{equation}
\label{eq:rel-7-162}
\phi_{\mathrm{rel}}(r)
  = (1-\mu)\Bigl(\frac{1}{r^{2}}-\kappa_{\mathrm{rel}}\Bigr)
    \gamma^{2}(r),
\end{equation}
integrate over the gap $\eta\le r\le 1$, and use the boundary
conditions (no-slip on both walls):
\begin{equation}
\label{eq:rel-7-163}
u=v=0, \qquad Du=0
\qquad\text{at } r=1 \text{ and } r=\eta.
\end{equation}
After integration by parts (using the same identities~\eqref{eq:7-172}
and~\eqref{eq:7-173} as in the classical case), one obtains an
equation of the form
\begin{equation}
\label{eq:rel-7-164}
\mathrm{re}(\sigma)\,
  \bigl\{I_{1}^{\mathrm{rel}}
    - \mathscr{T}_{\mathrm{rel}}\,a^{2}\,I_{3}^{\mathrm{rel}}
    + \mathrm{re}(I_{4}^{\mathrm{rel}})\bigr\} = 0,
\end{equation}
where the relativistic integral functionals are
\begin{align}
\label{eq:rel-7-165}
I_{1}^{\mathrm{rel}}
  &= \int_{\eta}^{1}\gamma^{2}(r)
     \Bigl\{r\Bigl|\frac{du}{dr}\Bigr|^{2}
     + \Bigl(\frac{1}{r}+a^{2}r\Bigr)|u|^{2}\Bigr\}\,dr,
\\[4pt]
I_{3}^{\mathrm{rel}}
  &= \int_{\eta}^{1}\phi_{\mathrm{rel}}(r)\,|v|^{2}\,dr,
\\[4pt]
\mathrm{re}(I_{4}^{\mathrm{rel}})
  &= \int_{\eta}^{1}r\,\gamma^{2}(r)\,\phi_{\mathrm{rel}}(r)
     \Bigl|\frac{dv}{dr}-\frac{v}{r}\Bigr|^{2}\,dr.
\end{align}
The factors $\gamma^{2}(r)\ge 1$ are everywhere positive and do not
alter the sign-definiteness of the integrands.  When $\mu>\eta^{2}$
we have $\mathscr{T}_{\mathrm{rel}}<0$ (just as in the classical
case, equation~\eqref{eq:7-189}), so the coefficient of
$\mathrm{re}(\sigma)$ in~\eqref{eq:rel-7-164} is
\emph{positive definite}.  Therefore
\begin{equation}
\label{eq:rel-7-166}
\boxed{\mathrm{re}(\sigma) < 0
  \qquad\text{for } \mu > \eta^{2},}
\end{equation}
and the relativistic viscous Couette flow is stable whenever
Rayleigh's criterion is satisfied.  This is the relativistic
generalisation of equation~\eqref{eq:7-189}.

\subsection{Eigenvalue problem for the relativistic Taylor
  number}\label{sec:rel-7-70e}

The marginal state is defined by $\mathrm{re}(p)=0$.  Following
the classical analysis, we expect the marginal state to be
\emph{stationary} ($p=0$, exchange of stabilities), so that
$\nu_{\mathrm{eff}}=\nu$ and the Israel--Stewart relaxation plays no
role at onset.  In this case
equations~\eqref{eq:rel-7-158}--\eqref{eq:rel-7-159} reduce to the
eigenvalue problem
\begin{equation}
\label{eq:rel-7-167}
\gamma^{2}(DD_{*}-a^{2})^{2}u
  = -T_{\mathrm{rel}}\,a^{2}
    \Bigl(\frac{1}{r^{2}}-\kappa_{\mathrm{rel}}\Bigr)v,
\end{equation}
\begin{equation}
\label{eq:rel-7-168}
\gamma^{2}(DD_{*}-a^{2})v = u,
\end{equation}
subject to the boundary conditions~\eqref{eq:rel-7-163}.  This is a
generalised eigenvalue problem for $T_{\mathrm{rel}}$ as a function
of the dimensionless wave number~$a$.  The critical relativistic
Taylor number is
\begin{equation}
\label{eq:rel-7-169}
T_{\mathrm{rel,cr}} = \min_{a}\,T_{\mathrm{rel}}(a)
  = T_{\mathrm{cl,cr}}
    \bigl[1+O(V^{2}/c^{2})\bigr],
\end{equation}
where $T_{\mathrm{cl,cr}}$ is the classical critical Taylor number.
The relativistic corrections are of order $V^{2}/c^{2}$ and
increase the critical Taylor number, i.e.\ relativity has a
\emph{stabilising} effect.  This is physically expected: the
enhanced inertia ($\gamma^{2}$ factors) makes it harder to displace
fluid elements against the centrifugal gradient.

\subsection{Role of Israel--Stewart relaxation for oscillatory
  modes}\label{sec:rel-7-70f}

While the Israel--Stewart relaxation drops out at marginal stability
($p=0$), it is essential for the \emph{dynamics} of perturbations
away from the marginal state.  For oscillatory modes ($p=i\omega$)
the effective viscosity~\eqref{eq:rel-7-157} becomes
\begin{equation}
\label{eq:rel-7-170}
\nu_{\mathrm{eff}}(i\omega)
  = \frac{\nu}{1+i\tau_{\pi}\omega}
  = \frac{\nu(1-i\tau_{\pi}\omega)}{1+\tau_{\pi}^{2}\omega^{2}},
\end{equation}
which has a reduced real part and a non-zero imaginary part.  The
damping rate of stable modes is therefore modified:
\begin{equation}
\label{eq:rel-7-171}
\mathrm{re}(p)
  = -\frac{\nu k^{2}}{1+\tau_{\pi}^{2}\omega^{2}}
    + \cdots,
\end{equation}
showing that high-frequency perturbations are
\emph{less efficiently damped} than in the Navier--Stokes theory.
This is a direct consequence of causality: in the Israel--Stewart
theory, viscous damping cannot act instantaneously.

More significantly, for any mode the viscous signal speed is bounded
by~\eqref{eq:rel-7-137}; perturbation fronts cannot outrun this
speed regardless of frequency.  In the Navier--Stokes limit
$\tau_{\pi}\to 0$ the signal speed diverges, violating causality.

\subsection{Summary of relativistic modifications}\label{sec:rel-7-70g}

The relativistic treatment of viscous Couette flow and its stability
introduces the following modifications to the classical theory of
\S\S\ref{sec:7-69}--\ref{sec:7-70}:

\begin{enumerate}
\item \textbf{Equilibrium profile:}  The velocity
  $V(r)=Ar+B/r$ receives corrections of order $V^{2}/c^{2}$
  (equation~\eqref{eq:rel-7-145}).
\item \textbf{Viscous constitutive law:}  The Navier--Stokes
  relation is replaced by the Israel--Stewart relaxation
  equation~\eqref{eq:rel-7-135}, ensuring that viscous signals
  propagate at finite speed $v_{\mathrm{visc}}\le c$.
\item \textbf{Perturbation equations:}  The eigenvalue problem
  acquires $\gamma^{2}$ factors and the frequency-dependent
  effective viscosity $\nu_{\mathrm{eff}}=\nu/(1+\tau_{\pi}p)$.
\item \textbf{Sufficient stability condition:}  The criterion
  $\mu>\eta^{2}$ remains sufficient for stability, with the proof
  carrying through thanks to the positive-definiteness of
  $\gamma^{2}$.
\item \textbf{Critical Taylor number:}  $T_{\mathrm{rel,cr}}>
  T_{\mathrm{cl,cr}}$; relativity stabilises the flow.
\item \textbf{Causality:}  All viscous propagation speeds are
  bounded by $v_{\mathrm{visc}}\le c$, guaranteed by
  $\tau_{\pi}>0$.
\end{enumerate}
